\documentclass[10pt,a4paper]{amsart}
\usepackage[margin=19mm]{geometry}
\usepackage{amsmath,amssymb,mathtools}
\usepackage[scr=rsfs,cal=euler]{mathalfa}
\usepackage{xfrac}
\usepackage[unicode,hidelinks,bookmarks=false]{hyperref}
\newcommand{\ie}{i.e.\ }
\newcommand{\cf}{cf.\ }
\newcommand{\resp}{respectively\ }
\newcommand{\Subsection}{\subsection}
\providecommand{\nobreakdash}{\nobreak\hbox{-}}
\setlength{\emergencystretch}{2em}
\multlinegap0pt
\usepackage{tikz}
\usepackage{amssymb}
\usepackage[shortlabels]{enumitem}
\usepackage{xcolor}
\usepackage[all]{xy}
\newtheorem{theorem}{Theorem}
\newtheorem{definition}{Definition}
\newcommand{\Aut}{\operatorname{Aut}}
\newcommand{\AutI}{\Aut_{\text{I}}}
\newcommand{\Autc}{\Aut_{\text{c}}}
\newcommand{\Autcb}{\Aut_{\text{cb}}}
\newcommand{\G}{\mathcal{G}}
\newcommand{\Ker}{\operatorname{Ker}}
\title{Refining incidence systems while preserving the automorphism pair}
\author{Antonio D\'iaz Ramos, R\'emi Molinier, Antonio Viruel}
\date{Source: arXiv:2511.02410v2 (CC BY 4.0)}
\begin{document}
\maketitle

\noindent\emph{Source and licence.} Every finite group is represented by a finite incidence geometry. Version arXiv:2511.02410v2 (CC BY 4.0), https://arxiv.org/abs/2511.02410v2. The exact-version arXiv abstract page was retrieved and its CC BY 4.0 article licence was independently checked on 2026-10-10. Licence and source notices: \texttt{licenses/arxiv-2511.02410-CC-BY-4.0.txt}.\par\smallskip

\noindent\textit{Editorial scope.} The counted unit is the original Theorem thm:improving incidence system (source0.tex lines 440--442), with its complete author proof at lines 443--585. The statement, proof and retained source passages are copied line for line without mathematical rewriting. Uncounted source-local context is retained at 181--203. The source article remains the work cited above; the title of this delivery describes the selected result. The unchanged native source is bundled as \texttt{sources/188412/source0.tex}. The original editable TikZ diagrams remain in the proof. The added definitions explain incidence, flags, correlations and the type-preserving automorphism subgroup. Difficulty is H2 after independent content review, replacing the earlier H3 label.
\begin{definition}
An \emph{incidence system} is a tuple $\Gamma=(X,\ast,t,I)$ with $X$ a non-empty set of \emph{elements}, $I$ a finite set of \emph{types}, $t\colon X\to I$ a surjective map called \emph{type function} and $\ast$ a symmetric binary relation called \emph{incidence relation}, such that for every $x,y\in X$, if $x\neq y$ and $x\ast y$, then $t(x)\neq t(y)$. The number of types $|I|$ is then the \emph{rank} of the incidence system. Finally for every $x\in X$ the \emph{degree} of $x$ is $|\{y\in X\mid\ y\ast x\}|$. 
\end{definition}

However, one way to think of an incidence system $\Gamma=(X,\ast,t,I)$ is as a simple proper colored graph $\G=(V,E,I,t)$ with set of vertices $V=X$, palette of colors $I$ and such that, for all $x,y\in V$, $\{x,y\}\in E$ if and only if $x\ast y$. This graph is commonly called \emph{incidence graph} of the incidence system $(X,\ast,t,I)$. In this paper, we will mostly work from the graph point-of-view. 
We now define the notion of incidence geometry. Notice also that, since an element cannot be incident to itself, the degree of an element $x\in X$ corresponds to the degree of $x$ seen as a vertex in the incidence graph.

\begin{definition}
Let $(X,\ast,t,I)$ be an incidence system. 
A \emph{flag} of an incidence system is a subset $F\subset X$ of pairwise incident elements, i.e. such that for every $x,y\in F$, if $x\neq y$ then $x\ast y$. We define the \emph{rank} of $F$ as its size,  $|F|$, or equivalently, as the number of types occurring in the flag. From the incidence graph point-of-view, a flag corresponds to a clique, and we define the \emph{rank} of a clique as its size. We define a \emph{chamber} as a flag of rank $|I|$ or as a clique of rank $|I|$. An incidence system is then an \emph{incidence geometry} if every flag is contained in a chamber or, equivalently, if every clique is contained in a chamber.
\end{definition}

Let us look at possible automorphisms of these incidence geometries.

\begin{definition}
Let $\Gamma=(X,\ast,t,I)$ be an incidence system. A \emph{correlation} of $\Gamma$ is a permutation $f$ of $X$ such that for all $x,y\in X$, 
\[ 
x\ast y\Leftrightarrow f(x)\ast f(y)\qquad \text{and}\qquad t(x)=t(y)\Leftrightarrow t(f(x))=t(f(y)).
\]
We denote by $\sigma_f\in \Sigma_I$ the permutation of types induced by $f$, and define an \emph{automorphism} of $\Gamma$ as a correlation $f$ with $\sigma_f$ equal to the identity permutation, i.e., such that $t(f(x))=t(x)$ for all $x\in X$. We denote by $\Aut(\Gamma)$ the group of correlations of $\Gamma$ and by $\AutI(\Gamma)$ the subgroup of $\Aut(\Gamma)$ consisting of the automorphisms of $\Gamma$, and we say that $\Gamma$ is an $(\Aut(\Gamma),\AutI(\Gamma))$-incidence system.
\end{definition}

From the graph point-of-view, a correlation of a incidence system $\Gamma$ is just an automorphism of the corresponding incidence graph $\G$ which is ``colored-blind", i.e., such that if two vertices have the same color, then their images have the same color.  Hence, we will call their graph counterpart \emph{colorblind automorphisms} and we will write $\Aut_{\text{cb}}(\G)$ for the set of these automorphisms of the incident graphs. Also, an automorphism of $\Gamma$ corresponds to an automorphism of $\G$ which preserves the colors. We call these automorphisms \emph{color automorphisms} and we write $\Autc(\G)$ for the subgroup consisting of these color automorphisms. So, if $\Gamma$ is an incidence system with incidence graph $\G$, we have $\Aut(\Gamma)=\Autcb(\G)$ and $\AutI(\Gamma)=\Autc(\G)$.

\begin{theorem}\label{thm:improving incidence system}
For any finite incidence system $\Gamma$, there exists a finite incidence system $\Gamma'$ with $(\Aut(\Gamma'),\AutI(\Gamma'))\cong (\Aut(\Gamma),\AutI(\Gamma))$ and such that all flags of $\Gamma'$ have rank at most two and all elements of $\Gamma'$ have degree at least two.
\end{theorem}

\begin{proof}
Let $\mathcal{G}=(V,E,I,t)$  be the incidence graph corresponding to $\Gamma$. Starting from $\mathcal{G}$, we will construct an incidence graph $\mathcal{G}'$ such that, if $\Gamma'$ is the incidence system associated to $\mathcal{G}'$, then $\Gamma'$ satisfies the conclusions of the statement.

Define $V_i$ as the subset of vertices of $V$ of degree $i$, set $V_{0,1}=V_0\cup V_1$, and consider the map $t\colon V\to I$ and the following natural number,
\[
M=\max_{i\in I}{|t^{-1}(\{i\})|},
\]
i.e., for every color $i\in I$, there is at most $M$ vertices of that same color. 

Next, we define a colored graph $\mathcal{G}'=(V',E',I',t')$ as follows,
\begin{align*}
V'&=V\sqcup \bigsqcup_{e\in E} \{u_{e,0},\ldots,u_{e,{2M+4}}\}\sqcup \bigsqcup_{w\in V_{0,1}} \{u_{w,0},\ldots,u_{w,{2M+4}}\},\\
E'&=\bigsqcup_{e=\{v,w\}\in E}\big\{\{v,u_{e,0}\},\{u_{e,0},w\},\{u_{e,0},u_{e,1}\},\{u_{e,1},u_{e,2}\},\ldots\\
&\ldots,\{u_{e,2M+3},u_{e,2M+4}\},\{u_{e,2M+4},u_{e,2M-1}\},\{u_{e,2M+3},u_{e,2M}\}\big\}\\
&\sqcup \bigsqcup_{w\in V_{0,1}} \big\{\{w,u_{w,0}\},\{u_{w,0},u_{w,1}\},\{u_{w,1},u_{w,2}\},\ldots\\
&\ldots,\{u_{w,2M+3},u_{w,2M+4}\},\{u_{w,2M+4},u_{w,2M-1}\},\{u_{w,2M+3},u_{w,2M}\}\big\},\\
&\sqcup \bigsqcup_{w\in V_0} \big\{\{w,u_{w,2}\}\big\},\\
I'&=I\sqcup \{i_0,i_1\},
\end{align*}
where the elements $u_{e,i}$ and $u_{v,i}$ are all different and do not belong to $V$, $i_0,i_1\notin I$, $t'(v)=t(v)$ for $v\in V$, and 
\[
t'(u_{e,j})=t'(u_{w,j})=\begin{cases} i_0&\text{if $j\equiv 0\pmod 2$},\\
i_1&\text{if $j\equiv 1\pmod 2$.}\\
\end{cases}
\]
Thus, we have replaced every edge $e=\{v,w\}$ of $E$ by its barycentric subdivision together with a ray of $2M$ vertices emanating from its midpoint, and with a figure eight involving $6$ vertices at the tail, and we have added to each vertex $w\in V_{0,1}$  a similar ray with and figure eight, plus an edge from $w$ to $u_{w,2}$ if $w\in V_0$. For instance, for $M=3$, we have, 
    
\begin{center}
\def\sl{0.2}%grid separation to control size of diagram
\begin{tikzpicture}[scale=5, font=\scriptsize]

%left-hand side vertical segment
\fill (0,\sl) circle (0.5pt);
\node[above=2pt] at (0,\sl) {$v$};
\draw (0,\sl)--(0,0)--(0,-\sl);
\fill (0,0) circle (0.5pt);
\node[left=2pt] at (0,0) {$u_{e,0}$};
\fill (0,-\sl) circle (0.5pt);
\node[below=2pt] at (0,-\sl) {$w$};

%central horizontal segment
\foreach \k in {1,2,3,4,5} {
    \fill (\k*\sl,0) circle (0.5pt);
    \node[below=2pt] at (\k*\sl,0) {$u_{e,\k}$};
    \draw (\k*\sl-\sl,0)--(\k*\sl,0);
}

%right-hand side figure eight
\fill (6*\sl,\sl) circle (0.5pt);
\node[above=2pt] at (6*\sl,\sl) {$u_{e,6}$};
\draw (5*\sl,0)--(6*\sl,\sl);
\fill (7*\sl,\sl) circle (0.5pt);
\node[above=2pt] at (7*\sl,\sl) {$u_{e,7}$};
\draw (6*\sl,\sl)--(7*\sl,\sl);
\fill (8*\sl,0) circle (0.5pt);
\node[right=2pt] at (8*\sl,0) {$u_{e,8}$};
\draw (7*\sl,\sl)--(8*\sl,0);
\fill (7*\sl,-\sl) circle (0.5pt);
\node[below=2pt] at (7*\sl,-\sl) {$u_{e,9}$};
\draw (8*\sl,0)--(7*\sl,-\sl);
\fill (6*\sl,-\sl) circle (0.5pt);
\node[below=2pt] at (6*\sl,-\sl) {$u_{e,10}$};
\draw (7*\sl,-\sl)--(6*\sl,-\sl)--(5*\sl,0);

\draw (7*\sl,-\sl)--(6*\sl,\sl);
  
\end{tikzpicture}
\end{center}
and
\begin{center}
\def\sl{0.2}%grid separation to control size of diagram
\begin{tikzpicture}[scale=5, font=\scriptsize]

%two left-most verteices
\draw[dashed,thick] (-2*\sl,0)--(-\sl,0);
\fill (-\sl,0) circle (0.5pt);
\node[below=2pt] at (-\sl,0) {$w$};
\draw (-\sl,0)--(0,0);
\draw[dotted,thick] (-\sl,0) arc[start angle=180, end angle=0, radius=1.5*\sl];
\fill (0,0) circle (0.5pt);
\node[below=2pt] at (0,0) {$u_{w,0}$};


%central horizontal segment
\foreach \k in {1,2,3,4,5} {
    \fill (\k*\sl,0) circle (0.5pt);
    \node[below=2pt] at (\k*\sl,0) {$u_{w,\k}$};
    \draw (\k*\sl-\sl,0)--(\k*\sl,0);
}

%right-hand side figure eight
\fill (6*\sl,\sl) circle (0.5pt);
\node[above=2pt] at (6*\sl,\sl) {$u_{w,6}$};
\draw (5*\sl,0)--(6*\sl,\sl);
\fill (7*\sl,\sl) circle (0.5pt);
\node[above=2pt] at (7*\sl,\sl) {$u_{w,7}$};
\draw (6*\sl,\sl)--(7*\sl,\sl);
\fill (8*\sl,0) circle (0.5pt);
\node[right=2pt] at (8*\sl,0) {$u_{w,8}$};
\draw (7*\sl,\sl)--(8*\sl,0);
\fill (7*\sl,-\sl) circle (0.5pt);
\node[below=2pt] at (7*\sl,-\sl) {$u_{w,9}$};
\draw (8*\sl,0)--(7*\sl,-\sl);
\fill (6*\sl,-\sl) circle (0.5pt);
\node[below=2pt] at (6*\sl,-\sl) {$u_{w,10}$};
\draw (7*\sl,-\sl)--(6*\sl,-\sl)--(5*\sl,0);

\draw (7*\sl,-\sl)--(6*\sl,\sl);
  
\end{tikzpicture}
\end{center}
where the dashed edge is present if and only if $w\in V_1$ and the dotted edge is present if and only if $w\in V_0$. Note also that
\begin{equation}\label{equ:large_t'_preimage}
 |t'^{-1}(i)|=\begin{cases} |t^{-1}(i)|\leq M &\text{if $i\neq i_0,i_1$},\\
(|E|+|V_{0,1}|)(M+3)&\text{if $i=i_0$},\\
(|E|+|V_{0,1}|)(M+2)&\text{if $i=i_1$}.
\end{cases}
\end{equation}
Moreover, as an incidence system has a non-empty set of elements by definition, we have $|E|+|V_{0,1}|>0$, $1\leq M<|t'^{-1}(i_0)|,|t'^{-1}(i_1)|$, and $|t'^{-1}(i_0)|\neq |t'^{-1}(i_1)|$. We also note that $\mathcal{G}'=(V',E',I',t')$ is a simple proper colored graph and that, by construction, $\mathcal{G}'$ has no clique of rank larger than two and has all elements of degree at least $2$.

Consider now a colorblind automorphism $f$ of $\mathcal{G}'$. Notice that, by \eqref{equ:large_t'_preimage}, the colors $i_0$ and $i_1$ are fixed. Thus, $f(V'\setminus V)=V'\setminus V$ and $f(V)=V$. Hence $f$ induces a colorblinded automorphism $f|_V$ of $\mathcal{G}$ and we have a group homomorphism
\[\begin{array}{ccccc}
\Phi & : & \Autcb(\G') & \to & \Autcb(\G) \\
 & & f & \mapsto & f|_V. \\
\end{array}\]

Let $f\in \Ker \Phi$. Since $f|_V$ is the identity and $f$ fixes the color $i_0$ then, for every $e\in E$, we must have $f(u_{e,0})=u_{e,0}$ and, for every $w\in V_{0,1}$, we must have $f(u_{w,0})=u_{w,0}$. As $f$ also fixes the color $i_1$ and the graph consisting of each ray and its corresponding figure eight has no non-trivial color automorphism, we also must have $f(u_{e,j})=u_{e,j}$ and $f(u_{w,j})=u_{w,j}$ for all $j\in\{0,1,\dots,2M+4\}$. Hence $f$ is the identity and we conclude that $\Phi$ is injective. 

Now let $f\in\Autcb(\mathcal{G})$ and consider $\widetilde{f}\colon V'\to V'$ defined, for all $v\in V$, by $\widetilde{f}(v)=f(v)$, and for all $e\in E$, $w\in V_{0,1}$, and $j\in\{0,1,\dots,2M+4\}$, by $\widetilde{f}(u_{e,j})=u_{f(e),j}$ and $\widetilde{f}(u_{w,j})=u_{f(w),j}$. Then, by construction, $\widetilde{f}$ is a graph automorphism. Moreover, if $\sigma_f\in \Sigma_I$ is the permutation of the colors $I$ induced by $f$, then $\widetilde{f}$ induces the permutation of colors $\sigma_{\widetilde{f}}\in \Sigma_{I'}$ that extends $\sigma_f$ by $i_j\mapsto i_j$ for $j=0,1$. Hence $\widetilde{f}\in\Autcb(\mathcal{G}')$ and, by construction, $\Phi(\widetilde{f})=f$. So $\Phi$ is an isomorphism. 

Finally, note that, for $f\in \Autcb(\G')$, we have that $\sigma_{\Phi(f)}$ is the restriction to $I$ of $\sigma_f$ and that $\sigma_f(i_j)=i_j$ for $j=0,1$. Hence, $\sigma_f$ is the identity permutation on $I'$ if and only if $\sigma_{\Phi(f)}$ is the identity permutation on $I$. Therefore $\Phi(\Autc(\Gamma'))=\Autc(\Gamma)$.
% _________________________________________

% Next, by \eqref{equ:large_t'_preimage} and the definition of $M$, every colorblind or color automorphism $f$ of $\Gamma'$ must send type $i_j$ to type $i_j$ for $j=0,1,2$. In addition, as the $u_{e,0}$'s with $e\in E$ are the only vertices of type $i_0$ with degree $3$, they must be permuted among themselves. We deduce that, in fact, we have 
% \begin{equation}\label{eqref:gamma_on_raytail}
% f(u_{e,j})=u_{f(e),j}
% \end{equation}
% for all $e\in E$ and all $0\leq j\leq 3M+5$, where, for $\{v,w\}\in E$, we define $f(e)=\{f(v),f(w)\}$. In the opposite direction, every colorblind or color automorphism of $\Gamma$ can be extended in a unique way to $\Gamma'$ using the formula \eqref{eqref:gamma_on_raytail}. Thus we have
% \[
% (\Autcb(\Gamma),\Autc(\Gamma))\cong (\Autcb(\Gamma'),\Autc(\Gamma'))
% \]
% and we are done.
\end{proof}

\end{document}
